\paragraph{末位终板刚性、临界预言机振荡与 Riesz 能量常数}
在末位相位的最终严格分层、bin-fold 规范群的最终满秩化、critical-scale 预言机容量骨架、escort 温度族的双折点响应以及六倍窗 exact dark band 的半阶指数尺度已经确立之后，还可把终板几何、商结构障碍、临界预算振荡与二阶能量常数进一步压缩为同一条后续链。在本段中统一记
\[
d_m(w):=d_m^{\mathrm{bin}}(w),
\qquad
\lambda:=\frac{2}{\varphi},
\qquad
\Omega_m:=\{0,1\}^m.
\]
下列结果分别刻画最小扇区的末位实现、大窗口上的全局非商化刚性、异常签名维数的 Fibonacci 饱和、临界比特率下的不可压平簇集带、escort 温度族的临界响应带、中心纤维场的精确二阶常数、碰撞概率与非平凡 Fourier 能量的精确主项，以及最大纤维超额的显式主系数。

\begin{theorem}[最小扇区的末位实现与终板逼近]\label{thm:xi-time-part70d-minsector-terminal-sheet-realization}
定义
\[
d_{\min}(m):=\min_{v\in X_m}d_m(v),
\qquad
A_1^{(m)}:=\{w\in X_m:\ w_m=1\},
\qquad
S_{m,\min}:=\{w\in X_m:\ d_m(w)=d_{\min}(m)\}.
\]
若“子覆盖同构”恒等式
\[
|S_{m,\min}|=|X_{m-2}|
\]
沿某个趋于无穷的偶数序列成立，则沿同一序列有
\[
\frac{|S_{m,\min}\triangle A_1^{(m)}|}{|X_m|}\longrightarrow 0.
\]
\end{theorem}
\leanverified{paper\_xi\_time\_part70d\_minsector\_terminal\_sheet\_realization}
\begin{proof}
由定理 \ref{thm:xi-time-part70ab-terminal-bit-eventual-strict-stratification}，对充分大的 \(m\) 有
\[
S_{m,\min}\subseteq A_1^{(m)}.
\]
另一方面，在假设下
\[
\frac{|S_{m,\min}|}{|X_m|}
=
\frac{|X_{m-2}|}{|X_m|}
=
\frac{F_m}{F_{m+2}}
\longrightarrow
\varphi^{-2},
\]
这也正是推论 \ref{cor:fold-bin-minsector-density-phi-minus2} 的内容。又由稳定词计数，
\[
|A_1^{(m)}|=F_m,
\qquad
|X_m|=F_{m+2},
\]
故
\[
\frac{|A_1^{(m)}|}{|X_m|}=\frac{F_m}{F_{m+2}}\longrightarrow \varphi^{-2}.
\]
由于 \(S_{m,\min}\subseteq A_1^{(m)}\) 且两者相对密度极限相同，遂有
\[
\frac{|A_1^{(m)}\setminus S_{m,\min}|}{|X_m|}\longrightarrow 0.
\]
而此时
\[
S_{m,\min}\triangle A_1^{(m)}=A_1^{(m)}\setminus S_{m,\min},
\]
故结论成立。
\end{proof}

\begin{theorem}[大窗口上的全局非商化刚性]\label{thm:xi-time-part70d-global-nonquotient-rigidity}
存在整数 \(m_0\)，使得对一切 \(m\ge m_0\)，等价关系
\[
\omega\sim_m\eta
\qquad\Longleftrightarrow\qquad
\Fold_m^{\mathrm{bin}}(\omega)=\Fold_m^{\mathrm{bin}}(\eta)
\]
既不能由任何有限群在 \(\Omega_m\) 上的自由作用之轨道划分实现，也不能来自任何有限阿贝尔群上的线性陪集划分。更准确地，不存在有限阿贝尔群 \(A_m\)、子群 \(H_m\le A_m\) 及双射
\[
\iota_m:\Omega_m\longrightarrow A_m
\]
使得
\[
\omega\sim_m\eta
\qquad\Longleftrightarrow\qquad
\iota_m(\omega)-\iota_m(\eta)\in H_m.
\]
\end{theorem}
\leanverified{paper\_xi\_time\_part70d\_global\_nonquotient\_rigidity}
\begin{proof}
记
\[
A_0^{(m)}:=\{w\in X_m:\ w_m=0\},
\qquad
A_1^{(m)}:=\{w\in X_m:\ w_m=1\}.
\]
由定理 \ref{thm:xi-time-part70ab-terminal-bit-eventual-strict-stratification}，存在 \(m_0\) 使得当 \(m\ge m_0\) 时，
\[
\min_{w\in A_0^{(m)}}d_m(w)>\max_{v\in A_1^{(m)}}d_m(v).
\]
由于 \(A_0^{(m)}\) 与 \(A_1^{(m)}\) 都非空，故对每个 \(m\ge m_0\)，纤维大小在同一窗口内至少取两个不同整数值。

若 \(\sim_m\) 来自某个有限群 \(G_m\) 的自由作用，则每个轨道大小都恒等于 \(|G_m|\)，因而所有纤维必同势；矛盾。

若 \(\sim_m\) 来自上式所述的线性陪集划分，则每个等价类都是某个陪集的原像，其大小恒等于 \(|H_m|\)，故所有纤维同样必同势；再度矛盾。于是两类实现均不可能存在。
\end{proof}

\begin{theorem}[异常签名维数的渐近饱和]\label{thm:xi-time-part70d-anomaly-signature-saturation}
\leanverified{paper\_xi\_time\_part70d\_anomaly\_signature\_saturation}
记
\[
G_m^{\mathrm{bin}}:=\prod_{w\in X_m}\mathfrak S_{d_m(w)}.
\]
则存在整数 \(m_0\)，使得对一切 \(m\ge m_0\)，都有
\[
\bigl(G_m^{\mathrm{bin}}\bigr)^{\mathrm{ab}}
\cong
(\ZZ/2\ZZ)^{|X_m|}
\cong
(\ZZ/2\ZZ)^{F_{m+2}}.
\]
因而任何保持“纤维独立置换”语义、并在 \(\FF_2\) 上完备分离 \(\bigl(G_m^{\mathrm{bin}}\bigr)^{\mathrm{ab}}\) 全部坐标的异常见证系统，其最小比特数恰为
\[
F_{m+2}.
\]
\end{theorem}
\begin{proof}
这正是定理 \ref{thm:xi-time-part65-binfold-center-vanishing-signature-stability} 的重述。该定理已证明：对充分大的 \(m\)，有
\[
\bigl(G_m^{\mathrm{bin}}\bigr)^{\mathrm{ab}}
\cong
(\ZZ/2\ZZ)^{F_{m+2}},
\qquad
\operatorname{rank}_{\FF_2}\!\bigl((G_m^{\mathrm{bin}})^{\mathrm{ab}}\bigr)=F_{m+2},
\]
并且纤维奇偶荷的最小完备秩在同一范围内严格等于 \(F_{m+2}\)。这恰好给出所述异常签名维数结论。
\end{proof}

\begin{theorem}[均匀微态输入下的临界比特率相变与簇集区间]\label{thm:xi-time-part70d-critical-bitrate-phase-cluster}
记
\[
\alpha:=\log_2\!\left(\frac{2}{\varphi}\right),
\qquad
\mathrm{Succ}_m(B):=\mathrm{Succ}_m^{\mathrm{bin}}(B).
\]
则成立：
\begin{enumerate}
  \item 若 \(\rho<\alpha\)，则
  \[
  \mathrm{Succ}_m(\lfloor \rho m\rfloor)\longrightarrow 0.
  \]
  \item 若 \(\rho>\alpha\)，则
  \[
  \mathrm{Succ}_m(\lfloor \rho m\rfloor)\longrightarrow 1.
  \]
  \item 对任意固定 \(\beta\in\RR\)，令
  \[
  B_m:=\lfloor \alpha m+\beta\rfloor.
  \]
  则序列 \(\bigl(\mathrm{Succ}_m(B_m)\bigr)_{m\ge 1}\) 不收敛，其全部簇集点恰组成闭区间
  \[
  \left[\frac{\varphi^2}{2\sqrt5},\,1\right].
  \]
\end{enumerate}
\end{theorem}
\begin{proof}
记 \(\lambda:=2/\varphi\)，则 \(\alpha=\log_2\lambda\)。

若 \(\rho<\alpha\)，则由平凡上界
\[
C_m^{\mathrm{bin}}(B)\le 2^B|X_m|
\]
得到
\[
\mathrm{Succ}_m(\lfloor \rho m\rfloor)
\le
\frac{2^{\lfloor \rho m\rfloor}F_{m+2}}{2^m}
=
\frac{2^{\lfloor \rho m\rfloor}}{\lambda^m}\cdot \frac{F_{m+2}}{\varphi^m}
\longrightarrow 0,
\]
因为 \(2^{\lfloor \rho m\rfloor}/\lambda^m\to 0\) 且 \(F_{m+2}/\varphi^m\to \varphi^2/\sqrt5\)。

若 \(\rho>\alpha\)，则 \(2^{\lfloor \rho m\rfloor}/\lambda^m\to+\infty\)。由推论 \ref{cor:xi-time-part70ab-extremal-fibers-terminal-bit-localization}，
\[
M_m:=\max_{w\in X_m}d_m(w)=\lambda^m+O(1).
\]
故对充分大的 \(m\) 有 \(2^{\lfloor \rho m\rfloor}>M_m\)，从而所有纤维都完全通过，即
\[
\mathrm{Succ}_m(\lfloor \rho m\rfloor)=1.
\]

现取临界序列 \(B_m=\lfloor \alpha m+\beta\rfloor\)。定义
\[
\theta_m:=\alpha m+\beta-\lfloor \alpha m+\beta\rfloor\in[0,1),
\qquad
\tau_m:=\frac{2^{B_m}}{\lambda^m}=2^{-\theta_m}\in[1/2,1].
\]
若 \(\log_2\varphi\in\QQ\)，则存在整数 \(p,q\ge 1\) 使 \(\varphi^q=2^p\in\QQ\)，这与 \(\varphi^q\notin\QQ\) 矛盾。故 \(\alpha=1-\log_2\varphi\) 为无理数，从而 \((\theta_m)\) 在 \([0,1]\) 上稠密，\((\tau_m)\) 的簇集点恰为整个区间 \([1/2,1]\)。

任取 \(\tau\in[1/2,1]\) 并取子列 \(m_j\) 使 \(\tau_{m_j}\to\tau\)。由定理 \ref{thm:xi-foldbin-dyadic-capacity-critical-window-limit}，
\[
\mathrm{Succ}_{m_j}(B_{m_j})
\longrightarrow
\Psi(\tau),
\]
其中
\[
\Psi(\tau)=
\begin{cases}
\dfrac{\varphi^2}{\sqrt5}\,\tau, & 0\le \tau\le \varphi^{-1},\\[6pt]
\dfrac{\varphi}{\sqrt5}\,\tau+\dfrac{1}{\varphi\sqrt5}, & \varphi^{-1}\le \tau\le 1,\\[8pt]
1, & \tau\ge 1.
\end{cases}
\]
因此 \(\bigl(\mathrm{Succ}_m(B_m)\bigr)\) 的全部簇集点恰为连续单调函数 \(\Psi\) 在 \([1/2,1]\) 上的像，即
\[
\Psi([1/2,1])
=
\left[\Psi(1/2),\,\Psi(1)\right]
=
\left[\frac{\varphi^2}{2\sqrt5},\,1\right].
\]
该区间非单点，故序列不收敛。
\end{proof}

\begin{theorem}[escort 温度族下的临界响应带]\label{thm:xi-time-part70d-escort-critical-response-band}
对任意 \(q\ge 0\)，定义 escort 加权的阈值响应
\[
\mathrm{Succ}_{m}^{(q)}(B)
:=
\sum_{w\in X_m}\pi_{m,q}(w)\min\!\left(1,\frac{2^B}{d_m(w)}\right),
\qquad
\theta(q):=\frac{1}{1+\varphi^{q+1}}.
\]
则同一临界比特率
\[
\alpha=\log_2\!\left(\frac{2}{\varphi}\right)
\]
给出 sharp transition：
\begin{enumerate}
  \item 若 \(\rho<\alpha\)，则
  \[
  \mathrm{Succ}_{m}^{(q)}(\lfloor \rho m\rfloor)\longrightarrow 0.
  \]
  \item 若 \(\rho>\alpha\)，则
  \[
  \mathrm{Succ}_{m}^{(q)}(\lfloor \rho m\rfloor)\longrightarrow 1.
  \]
  \item 若对固定 \(\beta\in\RR\) 令
  \[
  B_m=\lfloor \alpha m+\beta\rfloor,
  \]
  则序列 \(\bigl(\mathrm{Succ}_{m}^{(q)}(B_m)\bigr)_{m\ge 1}\) 不收敛，其全部簇集点恰为
  \[
  \left[\frac{1+(\varphi-1)\theta(q)}{2},\,1\right].
  \]
\end{enumerate}
\end{theorem}
\begin{proof}
若 \(\rho<\alpha\)，则 \(2^{\lfloor \rho m\rfloor}/\lambda^m\to 0\)。由定理 \ref{thm:fold-bin-two-state-asymptotic}，存在常数 \(C>0\) 使对充分大的 \(m\) 与全部 \(w\in X_m\) 都有
\[
d_m(w)\ge \lambda^m\left(\varphi^{-1}-C\Bigl(\frac{\varphi}{2}\Bigr)^m\right).
\]
于是
\[
\mathrm{Succ}_{m}^{(q)}(\lfloor \rho m\rfloor)
\le
\frac{2^{\lfloor \rho m\rfloor}}{\lambda^m\left(\varphi^{-1}-C(\varphi/2)^m\right)}
\longrightarrow 0.
\]

若 \(\rho>\alpha\)，则记
\[
M_m:=\max_{w\in X_m}d_m(w).
\]
由推论 \ref{cor:xi-time-part70ab-extremal-fibers-terminal-bit-localization}，\(M_m=\lambda^m+O(1)\)。故对充分大的 \(m\) 有 \(2^{\lfloor \rho m\rfloor}>M_m\)，于是每个纤维都完全通过，从而
\[
\mathrm{Succ}_{m}^{(q)}(\lfloor \rho m\rfloor)=1.
\]

现令 \(B_m=\lfloor \alpha m+\beta\rfloor\)，并仍记
\[
\tau_m:=\frac{2^{B_m}}{\lambda^m}\in[1/2,1].
\]
由上一条定理的证明，\((\tau_m)\) 的全部簇集点恰为 \([1/2,1]\)。

记
\[
I:=[1/2,1],
\qquad
J:=\left[\frac{\varphi^{-1}}{2},\,2\right].
\]
函数
\[
f(\tau,x):=\min\!\left(1,\frac{\tau}{x}\right)
\]
在紧集 \(I\times J\) 上联合 Lipschitz。再由定理 \ref{thm:fold-bin-two-state-asymptotic}，
\[
\left|\frac{d_m(w)}{\lambda^m}-\varphi^{-w_m}\right|
\le
C\Bigl(\frac{\varphi}{2}\Bigr)^m
\qquad (w\in X_m),
\]
故对充分大的 \(m\) 与全部 \(w\in X_m\) 都有
\[
\min\!\left(1,\frac{2^{B_m}}{d_m(w)}\right)
=
f\!\left(\tau_m,\frac{d_m(w)}{\lambda^m}\right)
=
f\!\left(\tau_m,\varphi^{-w_m}\right)
+O\!\left(\left(\frac{\varphi}{2}\right)^m\right),
\]
且误差对 \(w\) 与 \(\tau_m\in I\) 一致。对 \(\pi_{m,q}\) 取期望，并用命题 \ref{prop:fold-bin-escort-lastbit}
\[
\mathbf P_{\pi_{m,q}}(w_m=1)=\theta(q)+O\!\left(\left(\frac{\varphi}{2}\right)^m\right),
\qquad
\mathbf P_{\pi_{m,q}}(w_m=0)=1-\theta(q)+O\!\left(\left(\frac{\varphi}{2}\right)^m\right),
\]
得到
\[
\mathrm{Succ}_{m}^{(q)}(B_m)
=
\bigl(1-\theta(q)\bigr)\min\{1,\tau_m\}
+
\theta(q)\min\{1,\varphi\tau_m\}
+o(1).
\]
任取 \(\tau\in[1/2,1]\) 及子列 \(m_j\) 使 \(\tau_{m_j}\to\tau\)，便有
\[
\mathrm{Succ}_{m_j}^{(q)}(B_{m_j})
\longrightarrow
\Psi_q^{\mathrm{esc}}(\tau),
\qquad
\Psi_q^{\mathrm{esc}}(\tau)
:=
\bigl(1-\theta(q)\bigr)\min\{1,\tau\}
+
\theta(q)\min\{1,\varphi\tau\}.
\]
故全部簇集点恰为 \(\Psi_q^{\mathrm{esc}}([1/2,1])\)。由于 \(\Psi_q^{\mathrm{esc}}\) 在 \([1/2,1]\) 上连续单调，故该像为区间
\[
\left[\Psi_q^{\mathrm{esc}}(1/2),\,\Psi_q^{\mathrm{esc}}(1)\right].
\]
又因 \(1/2<\varphi^{-1}\)，故
\[
\Psi_q^{\mathrm{esc}}(1/2)
=
\frac{1-\theta(q)}{2}+\frac{\varphi\theta(q)}{2}
=
\frac{1+(\varphi-1)\theta(q)}{2},
\qquad
\Psi_q^{\mathrm{esc}}(1)=1.
\]
于是所求簇集区间即为
\[
\left[\frac{1+(\varphi-1)\theta(q)}{2},\,1\right].
\]
其长度严格为正，故序列不收敛。
\end{proof}

\begin{theorem}[中心纤维场的精确二阶常数]\label{thm:xi-time-part70d-central-fiber-second-moment}
\leanverified{paper\_xi\_time\_part70d\_central\_fiber\_second\_moment}
令 \(U_m\sim \mathrm{Unif}(X_m)\)，并定义尺度化纤维场
\[
Y_m:=\left(\frac{\varphi}{2}\right)^m d_m(U_m).
\]
则有极限
\[
\mathbb E(Y_m^2)\longrightarrow 2\varphi^{-2},
\qquad
\operatorname{Var}(Y_m)\longrightarrow \varphi^{-7}.
\]
\end{theorem}
\begin{proof}
由推论 \ref{cor:xi-time-part70a-scaled-moments-leading-constant}，当 \(q=1,2\) 时分别有
\[
\mathbb E(Y_m)
\longrightarrow
\frac{1}{\varphi}+\varphi^{-3},
\qquad
\mathbb E(Y_m^2)
\longrightarrow
\frac{1}{\varphi}+\varphi^{-4}.
\]
利用恒等式 \(\varphi^2=\varphi+1\) 与 \(\sqrt5=2\varphi-1\)，可将上式化简为
\[
\frac{1}{\varphi}+\varphi^{-4}=2\varphi^{-2},
\qquad
\frac{1}{\varphi}+\varphi^{-3}=\frac{\sqrt5}{\varphi^2}.
\]
因此
\[
\operatorname{Var}(Y_m)
=
\mathbb E(Y_m^2)-\bigl(\mathbb E(Y_m)\bigr)^2
\longrightarrow
2\varphi^{-2}-\frac{5}{\varphi^4}.
\]
再用 \(\varphi^2=\varphi+1\) 化简，得到
\[
2\varphi^{-2}-\frac{5}{\varphi^4}=\varphi^{-7}.
\]
证毕。
\end{proof}

\begin{theorem}[碰撞概率与非平凡 Fourier 能量的精确主项]\label{thm:xi-time-part70d-collision-riesz-main-terms}
沿用定理 \ref{thm:xi-time-part70d-central-fiber-second-moment} 的记号，并沿用定理 \ref{thm:xi-time-part9p-chi2-ranktwo-resonant-floor} 的 Fourier 记号。记
\[
\mathsf{Col}_m:=
\frac{1}{4^m}\sum_{w\in X_m}d_m(w)^2.
\]
则有
\[
\mathsf{Col}_m\sim \frac{2}{\sqrt5}\,\varphi^{-m},
\qquad
F_{m+2}\,\mathsf{Col}_m\longrightarrow \frac{2\varphi^2}{5}.
\]
进一步，非平凡 Fourier 能量满足
\[
4^{-m}\sum_{k\ne 0}\bigl|\widehat d_m(k)\bigr|^2
\longrightarrow
\frac{1}{5\varphi^3}.
\]
\end{theorem}
\begin{proof}
由 \(U_m\sim\mathrm{Unif}(X_m)\) 的定义，
\[
\mathsf{Col}_m
=
\frac{|X_m|}{4^m}\,\mathbb E\!\bigl(d_m(U_m)^2\bigr).
\]
又因为
\[
d_m(U_m)=\lambda^m Y_m,
\qquad
\lambda=\frac{2}{\varphi},
\]
故
\[
\mathsf{Col}_m
=
|X_m|\,\varphi^{-2m}\,\mathbb E(Y_m^2).
\]
由 \(|X_m|=F_{m+2}\) 与定理 \ref{thm:xi-time-part70d-central-fiber-second-moment}，
\[
\mathsf{Col}_m
=
F_{m+2}\varphi^{-2m}\bigl(2\varphi^{-2}+o(1)\bigr).
\]
再用 Binet 公式
\[
F_{m+2}\sim \frac{\varphi^{m+2}}{\sqrt5},
\]
便得
\[
\mathsf{Col}_m
\sim
\frac{\varphi^{m+2}}{\sqrt5}\cdot \varphi^{-2m}\cdot 2\varphi^{-2}
=
\frac{2}{\sqrt5}\,\varphi^{-m}.
\]
两边再乘以 \(F_{m+2}\) 即得
\[
F_{m+2}\,\mathsf{Col}_m\longrightarrow \frac{2\varphi^2}{5}.
\]

另一方面，由定理 \ref{thm:xi-time-part9p-chi2-ranktwo-resonant-floor} 的第一条精确恒等式，
\[
F_{m+2}\,\mathsf{Col}_m-1
=
4^{-m}\sum_{k\ne 0}\bigl|\widehat d_m(k)\bigr|^2.
\]
取极限并用
\[
\frac{2\varphi^2}{5}-1=\frac{1}{5\varphi^3},
\]
便得到
\[
4^{-m}\sum_{k\ne 0}\bigl|\widehat d_m(k)\bigr|^2
\longrightarrow
\frac{1}{5\varphi^3}.
\]
证毕。
\end{proof}

\begin{corollary}[最大纤维超额的精确主项]\label{cor:xi-time-part70d-maxfiber-excess-exact-main-term}
记
\[
\bar d_m:=\frac{2^m}{F_{m+2}},
\qquad
M_m:=\max_{w\in X_m}d_m(w).
\]
则有
\[
M_m-\bar d_m
=
\varphi^{-4}\left(\frac{2}{\varphi}\right)^m
+o\!\left(\left(\frac{2}{\varphi}\right)^m\right).
\]
\end{corollary}
\begin{proof}
由推论 \ref{cor:xi-time-part70ab-extremal-fibers-terminal-bit-localization}，
\[
M_m=\left(\frac{2}{\varphi}\right)^m+O(1).
\]
另一方面，由 Binet 公式，
\[
\bar d_m
=
\frac{2^m}{F_{m+2}}
=
\frac{\sqrt5}{\varphi^2}\left(\frac{2}{\varphi}\right)^m
+o\!\left(\left(\frac{2}{\varphi}\right)^m\right).
\]
由于 \(O(1)=o((2/\varphi)^m)\)，相减即得
\[
M_m-\bar d_m
=
\left(1-\frac{\sqrt5}{\varphi^2}\right)\left(\frac{2}{\varphi}\right)^m
+o\!\left(\left(\frac{2}{\varphi}\right)^m\right).
\]
再用 \(\sqrt5=2\varphi-1\) 与 \(\varphi-1=\varphi^{-1}\) 化简：
\[
1-\frac{\sqrt5}{\varphi^2}
=
1-\frac{2\varphi-1}{\varphi^2}
=
\frac{(\varphi-1)^2}{\varphi^2}
=
\varphi^{-4}.
\]
故结论成立。
\end{proof}

\noindent
由此，末位 \(1\) 终板不再只是最小纤维的容纳层，而在“子覆盖同构”外推下渐近地与整个最小扇区重合；与此同时，充分大窗口上的折叠关系也被彻底排除出自由商与线性陪集两类等势模型。再向临界资源侧推进，dyadic 预算与 escort 温度都在同一比特率 \(\log_2(2/\varphi)\) 处出现不可压平的簇集带，而其下边界由黄金比例与温度参数共同给出显式常数。最后，中心纤维场的二阶矩、碰撞概率、非平凡 Fourier 能量与最大纤维超额都锁定到闭式黄金常数，从而终板刚性、资源相变与频域能量三条先前分立的链被压缩为同一条可审计的 Fibonacci 后续主线。

\endinput
